\chapter{Foreign Function Interface}
\index{FFI}\index{foreign function interface}
\label{ch:ffi}

Sometimes you need to use a library written in C---a graphics toolkit, a database driver, a compression algorithm, or a system API. Kaappi's Foreign Function Interface (FFI) lets you call C functions directly from Scheme without writing any C code yourself. This chapter covers opening shared libraries, binding functions, handling types, and passing Scheme callbacks into C.

\section{The \texttt{(kaappi ffi)} Library}

The FFI is provided by the \texttt{(kaappi ffi)}\index{kaappi ffi} library:

\begin{lstlisting}
(import (kaappi ffi))
\end{lstlisting}

The core procedures are:

\begin{center}
\begin{tabular}{lp{6.5cm}}
\textbf{Procedure} & \textbf{Purpose} \\
\hline
\texttt{ffi-open} & Open a shared library (.dylib / .so) \\
\texttt{ffi-fn} & Bind a C function from the library \\
\texttt{ffi-close} & Close the library handle \\
\texttt{ffi-callback} & Create a C function pointer from a Scheme procedure \\
\texttt{ffi-callback-release} & Free a callback \\
\texttt{ffi-callback?} & Test if a value is an FFI callback \\
\texttt{ffi-bytevector-ptr} & Get the data pointer of a bytevector (for passing buffers to C) \\
\end{tabular}
\end{center}

\section{Opening a Shared Library}
\index{ffi-open}

\begin{lstlisting}[style=repl]
> (import (kaappi ffi))
> (define libm (ffi-open "libm.dylib"))  ; macOS
\end{lstlisting}

On Linux, use the \texttt{.so}\index{shared library} name:

\begin{lstlisting}
(define libm (ffi-open "libm.so.6"))
\end{lstlisting}

The library is loaded into the process and stays open until you call \texttt{ffi-close} or the program exits.

\begin{note}[Library names by platform]
macOS uses \texttt{.dylib} extensions (\texttt{libm.dylib}, \texttt{libsqlite3.dylib}). Linux uses \texttt{.so} with version suffixes (\texttt{libm.so.6}, \texttt{libsqlite3.so.0}). You can also pass an absolute path.
\end{note}

\section{Binding C Functions}
\index{ffi-fn}

\texttt{ffi-fn} looks up a symbol in the library and creates a callable Scheme procedure:

\begin{lstlisting}
(ffi-fn library "function-name" '(param-types ...) 'return-type)
\end{lstlisting}

A function can take at most 16 parameters.

\subsection{Example: Math Functions}

\begin{lstlisting}[style=repl]
> (define c-sqrt (ffi-fn libm "sqrt" '(double) 'double))
> (define c-pow (ffi-fn libm "pow" '(double double) 'double))
> (define c-ceil (ffi-fn libm "ceil" '(double) 'double))
> (c-sqrt 2.0)
1.4142135623730951
> (c-pow 2.0 10.0)
1024.0
> (c-ceil 3.2)
4.0
\end{lstlisting}

The bound function behaves like any other Scheme procedure---you can pass it to \texttt{map}, store it in a variable, or use it in higher-order patterns.

\subsection{Example: String Functions}

\begin{lstlisting}[style=repl]
> (define libc (ffi-open "libc.dylib"))
> (define c-strlen (ffi-fn libc "strlen" '(string) 'size_t))
> (c-strlen "hello")
5
> (c-strlen "")
0
\end{lstlisting}

\section{Supported Types}
\index{FFI types}

The following C types are supported in parameter and return type positions:

\begin{center}
\begin{tabular}{ll}
\textbf{FFI type} & \textbf{C type} \\
\hline
\texttt{int} & \texttt{int} (32-bit signed) \\
\texttt{long} & \texttt{long} (64-bit on 64-bit systems) \\
\texttt{int8} / \texttt{uint8} & \texttt{int8\_t} / \texttt{uint8\_t} \\
\texttt{int16} / \texttt{uint16} & \texttt{int16\_t} / \texttt{uint16\_t} \\
\texttt{int32} / \texttt{uint32} & \texttt{int32\_t} / \texttt{uint32\_t} \\
\texttt{int64} / \texttt{uint64} & \texttt{int64\_t} / \texttt{uint64\_t} \\
\texttt{size\_t} & \texttt{size\_t} \\
\texttt{float} & \texttt{float} (32-bit) \\
\texttt{double} & \texttt{double} (64-bit) \\
\texttt{char} & \texttt{char} \\
\texttt{bool} & \texttt{\_Bool} \\
\texttt{string} & \texttt{const char*} (null-terminated) \\
\texttt{pointer} & \texttt{void*} (opaque pointer) \\
\texttt{void} & \texttt{void} (return type only) \\
\end{tabular}
\end{center}

\subsection{Type Conversions}

\begin{itemize}
    \item \textbf{Strings:} Scheme strings are automatically converted to null-terminated C strings when passed as arguments, and C strings are converted back to Scheme strings when returned.
    \item \textbf{Numbers:} Scheme integers and floats are converted to the appropriate C numeric type. Out-of-range values are rejected with an error naming the FFI function and the expected type.
    \item \textbf{Characters:} The \texttt{char} type takes and returns Scheme characters---a C function declared \texttt{(char) -> char} maps \texttt{\#\textbackslash a} to \texttt{\#\textbackslash A}, not \texttt{97} to \texttt{65}.
    \item \textbf{Pointers:}\index{pointer} Opaque pointer values are represented as Scheme FFI pointer objects. You cannot dereference them directly from Scheme.
    \item \textbf{Void:} Use \texttt{'void} as the return type for C functions that return nothing.
\end{itemize}

\section{A Complete Example: System Information}

Here is a more realistic example that combines several libc functions to query system information:

\begin{lstlisting}
(import (kaappi ffi))

(define libc (ffi-open "libc.dylib"))

;; Bind functions we need
(define c-getenv
  (ffi-fn libc "getenv" '(string) 'string))
(define c-strlen
  (ffi-fn libc "strlen" '(string) 'size_t))
(define c-time
  (ffi-fn libc "time" '(pointer) 'long))
(define c-getpid
  (ffi-fn libc "getpid" '() 'int))
\end{lstlisting}

\begin{lstlisting}[style=repl]
> (c-getenv "HOME")
"/Users/alice"    (*@\emph{(your values will differ)}@*)
> (c-getenv "SHELL")
"/bin/zsh"
> (c-strlen (c-getenv "HOME"))
12
> (c-getpid)
42137
> (c-time 0)
1751234567
\end{lstlisting}

\begin{lstlisting}
;; Build a system report
(define (system-info)
  (list
    (cons "home" (c-getenv "HOME"))
    (cons "shell" (c-getenv "SHELL"))
    (cons "pid" (c-getpid))
    (cons "time" (c-time 0))))

;; Clean up
(ffi-close libc)
\end{lstlisting}

\begin{note}[Why not SQLite?]
The SQLite C API uses pointer-to-pointer out-parameters (\texttt{sqlite3**}) to return database handles, which requires low-level pointer manipulation beyond what the FFI provides directly. In practice, you would use the \texttt{kaappi-sqlite} package (\texttt{thottam install kaappi-sqlite}), which wraps these details in a clean Scheme interface.
\end{note}

\section{Passing Buffers with Bytevectors}
\index{ffi-bytevector-ptr}\index{bytevectors!and FFI}

Many C functions do not return a string---they fill a buffer you hand them. \texttt{ffi-bytevector-ptr} returns the address of a bytevector's data, which you can pass wherever C expects a pointer. Here \texttt{getcwd} writes the current directory into a bytevector as a 0-terminated C string:

\begin{lstlisting}
(import (kaappi ffi))

(define libc (ffi-open "libc.dylib"))
(define c-getcwd
  (ffi-fn libc "getcwd" '(pointer size_t) 'pointer))

;; Extract the C string that ends at the first 0 byte
(define (bytevector-cstring bv)
  (let loop ((i 0))
    (if (zero? (bytevector-u8-ref bv i))
        (utf8->string (bytevector-copy bv 0 i))
        (loop (+ i 1)))))

(define buf (make-bytevector 4096 0))
(define result (c-getcwd (ffi-bytevector-ptr buf) 4096))
(display (bytevector-cstring buf))
\end{lstlisting}

\begin{codeoutput}
/Users/alice/projects
\end{codeoutput}

Keep the bytevector reachable for as long as C holds the pointer. Kaappi's garbage collector does not move objects, so the address stays valid while the bytevector is alive---but if the bytevector is collected, the pointer dangles.

\section{Callbacks: Passing Scheme to C}
\index{ffi-callback}

Some C APIs expect function pointers as arguments (e.g., comparators for \texttt{qsort}, event handlers, iteration callbacks). \texttt{ffi-callback} creates a C-compatible function pointer from a Scheme procedure:

\begin{lstlisting}
(define cb (ffi-callback scheme-proc '(param-types ...) 'return-type))
\end{lstlisting}

Callbacks support only a fixed set of signatures:

\begin{center}
\begin{tabular}{ll}
\textbf{Parameters} & \textbf{Return type} \\
\hline
\texttt{()} & \texttt{void} \\
\texttt{(int)} & \texttt{void} \\
\texttt{(pointer)} & \texttt{void} \\
\texttt{(pointer)} & \texttt{int} \\
\texttt{(int pointer)} & \texttt{int} \\
\texttt{(pointer pointer)} & \texttt{void} \\
\texttt{(pointer pointer)} & \texttt{int} \\
\end{tabular}
\end{center}

Any other combination raises \texttt{unsupported callback signature}. At most 32 callbacks can exist at once, which is another reason to release them promptly.

\subsection{Example: Sorting Bytes with \texttt{qsort}}

Pointer arguments arrive in the Scheme procedure as raw addresses. You cannot dereference them from Scheme---but you can hand them straight back to C. This comparator sorts the bytes of a bytevector with C's \texttt{qsort} by delegating each comparison to \texttt{memcmp}:

\begin{lstlisting}
(define c-qsort
  (ffi-fn libc "qsort"
    '(pointer size_t size_t pointer) 'void))
(define c-memcmp
  (ffi-fn libc "memcmp" '(pointer pointer size_t) 'int))

;; qsort passes addresses of two elements; compare 1 byte each
(define byte-compare
  (ffi-callback
    (lambda (a b) (c-memcmp a b 1))
    '(pointer pointer)
    'int))

(define data (bytevector 3 1 4 1 5 9 2 6))
(c-qsort (ffi-bytevector-ptr data) 8 1 byte-compare)
(ffi-callback-release byte-compare)
(display data)
\end{lstlisting}

\begin{codeoutput}
\#u8(1 1 2 3 4 5 6 9)
\end{codeoutput}

\begin{warning}[Always release callbacks]
Callbacks pin Scheme closures so the garbage collector cannot reclaim them. Always call \texttt{ffi-callback-release}\index{ffi-callback-release} when the callback is no longer needed to prevent memory leaks.
\end{warning}

If the Scheme procedure raises an error while C is calling it, the error is not lost: it is re-raised in your program as soon as the C call that invoked the callback returns, where \texttt{guard} can catch it as usual. The C code itself never sees a Scheme exception---it runs to completion first.

\section{Closing Libraries}
\index{ffi-close}

When you are done with a library, close it to release the handle:

\begin{lstlisting}[style=repl]
> (ffi-close libm)
\end{lstlisting}

After closing, any functions bound from that library become invalid. In practice, most programs open libraries at startup and never close them (the OS reclaims everything on exit).

\section{Error Handling}

FFI calls can crash your program if you get the types wrong---there is no safety net between Scheme and C. Common pitfalls:

\begin{itemize}
    \item \textbf{Wrong types:} Passing an integer where a pointer is expected (or vice versa) causes undefined behavior.
    \item \textbf{Null pointers:} C functions may return NULL; check pointer values before using them.
    \item \textbf{Lifetime issues:} If you pass a Scheme string to C and the C code stores the pointer for later use, the Scheme GC may collect the string. Copy the data on the C side if it needs to persist.
    \item \textbf{Thread safety:} Do not call FFI functions from multiple Kaappi fibers simultaneously unless the C library is thread-safe.
\end{itemize}

\begin{warning}[FFI bypasses safety]
The FFI is inherently unsafe. Incorrect type declarations, buffer overflows, and use-after-free in C code can crash the Kaappi process. Test FFI bindings carefully and prefer ecosystem packages (\texttt{kaappi-sqlite}, \texttt{kaappi-pg}, etc.) when available.
\end{warning}

The FFI error messages you are most likely to encounter, with their fixes, are listed in Appendix~\ref{app:errors}.

\subsection{Debugging FFI Problems}

When FFI calls misbehave, use this checklist:

\begin{enumerate}
    \item \textbf{Verify the library loaded.} If \texttt{ffi-open} returns without error, the library is found. If it raises an error, check the file name and path. On macOS, use \texttt{otool -L} to inspect a binary's dependencies; on Linux, use \texttt{ldd}.
    \item \textbf{Check the symbol name.} C++ name-mangling, static linking, and version suffixes can cause \texttt{ffi-fn} to fail. Use \texttt{nm -gU libfoo.dylib | grep function\_name} to confirm the symbol exists.
    \item \textbf{Match types exactly.} A function declared as returning \texttt{int} but actually returning \texttt{size\_t} may work on some platforms and silently corrupt values on others. Consult the C header file.
    \item \textbf{Test in isolation.} Bind one function at a time in the REPL. Call it with known inputs and compare against expected outputs before integrating into a larger program.
\end{enumerate}

\section{FFI and Sandbox Mode}
\index{sandbox}\index{--sandbox}

Kaappi's \texttt{--sandbox} mode disables the FFI entirely. This is useful for running untrusted code:

\begin{lstlisting}
kaappi --sandbox untrusted.scm
\end{lstlisting}

In sandbox mode, \texttt{(import (kaappi ffi))} itself fails with \texttt{sandbox: cannot load library from file}, and none of the FFI procedures are defined. See Appendix~\ref{app:cli} for the full set of command-line options, including \texttt{--sandbox}.

\section{When to Use the FFI}

\begin{itemize}
    \item \textbf{Use it} when you need a C library that has no Scheme wrapper (niche libraries, system APIs, hardware interfaces).
    \item \textbf{Prefer ecosystem packages} (\texttt{thottam install ...}) when they exist---they handle error checking, memory management, and provide idiomatic Scheme interfaces.
    \item \textbf{Consider writing a library} if you find yourself wrapping more than a few functions. Publish it via thottam so others benefit.
\end{itemize}

\section*{Summary}

\begin{itemize}
    \item The \texttt{(kaappi ffi)} library calls C functions in shared libraries at runtime.
    \item \texttt{ffi-open} loads a shared library and \texttt{ffi-fn} binds a C function to a Scheme procedure given its type signature.
    \item The FFI supports integer, floating-point, string, and pointer types; bytevectors pass buffers via \texttt{ffi-bytevector-ptr}, and Scheme procedures become C function pointers via \texttt{ffi-callback} (fixed signature set, released with \texttt{ffi-callback-release}).
    \item The FFI is inherently unsafe and is disabled by \texttt{--sandbox}; prefer ecosystem packages when they exist.
\end{itemize}

\section*{Exercises}

\begin{exercise}[Binding \texttt{strlen}]
Use the FFI to bind \code{strlen} from the C standard library. Open the appropriate shared library for your platform (\code{libc.dylib} on macOS, \code{libc.so.6} on Linux), bind \code{strlen} with the correct type signature, and test it:

\begin{lstlisting}[style=repl]
> (c-strlen "hello")
5
> (c-strlen "")
0
> (c-strlen "Kaappi Scheme")
13
\end{lstlisting}

What happens if you declare the return type as \code{'int} instead of \code{'size\_t}? Does it still work for short strings? Think about why \code{size\_t} is the correct choice (hint: consider strings longer than $2^{31}$ bytes).
\end{exercise}

\begin{exercise}[Comparing FFI \texttt{pow} with \texttt{expt}]
Bind the C function \code{pow} from \code{libm} using the FFI and write a procedure that compares its results with Scheme's built-in \code{expt}:

\begin{lstlisting}
(define (compare-pow base exp)
  (let ((c-result (c-pow (exact->inexact base)
                         (exact->inexact exp)))
        (s-result (expt base exp)))
    (list c-result s-result (- c-result s-result))))
\end{lstlisting}

Test with several inputs: \code{(compare-pow 2 10)}, \code{(compare-pow 3.0 3.0)}, and \code{(compare-pow 2 0.5)}. When do the results differ? Why might the C version and the Scheme version give slightly different answers? (Hint: think about exact vs.\ inexact arithmetic.)
\end{exercise}

\begin{exercise}[Safe Library Wrapper]
Write a procedure \code{(with-library name proc)} that opens a shared library, passes the handle to \code{proc}, and guarantees that \code{ffi-close} is called even if \code{proc} raises an exception. This is the FFI equivalent of Python's \code{with open(...) as f:} pattern.

\begin{lstlisting}
(define (with-library name proc)
  ;; Your implementation here.
  ;; Use guard or dynamic-wind to ensure cleanup.
  ...)
\end{lstlisting}

\code{guard} (Chapter~\ref{ch:errors}) is sufficient here; \code{dynamic-wind} (Chapter~\ref{ch:continuations}) is the more general tool for paired acquire/release. Test it by binding and calling \code{ceil} from libm:

\begin{lstlisting}[style=repl]
> (with-library "libm.dylib"
    (lambda (lib)
      (let ((c-ceil (ffi-fn lib "ceil"
                      '(double) 'double)))
        (c-ceil 3.2))))
4.0
\end{lstlisting}

Verify that the library is closed after the call returns. What happens if you try to use a function bound from a closed library?
\end{exercise}

In the next chapter, we cover concurrency---fibers, channels, and threads.

\chapterend
